\chapter*{Conclusion}
\addcontentsline{toc}{chapter}{Conclusion}
\markboth{Conclusion}{Conclusion}

This dissertation addressed fundamental challenges in developing adversarially resilient artificial intelligence for hybrid intrusion detection and prevention systems through integrated theoretical frameworks and practical implementations. The research culminated in two complementary master problems decomposing into fourteen sub-problems addressed by seven novel methods validated through comprehensive experimental evaluation across multiple benchmark datasets. This concluding chapter synthesizes key contributions, discusses broader implications for network security practice, acknowledges limitations, and outlines promising directions for future research including postdoctoral investigations.

\section*{Summary of Contributions}

The dissertation made theoretical, algorithmic, and empirical contributions advancing the state of the art in adversarially resilient network security.

\subsection*{Theoretical Foundations}

We established unified mathematical frameworks integrating continuous-time graph neural ordinary differential equations, differentially private optimal transport, PAC-Bayesian generalization bounds, variational Bayesian inference, and game-theoretic adversarial modeling. The theoretical analysis provided convergence guarantees for continuous-time graph neural ODEs through Lipschitz continuity analysis, Byzantine-robust convergence for federated learning under malicious participants achieving convergence despite adversarial interference, PAC-Bayesian generalization bounds for state space models enabling certified confidence intervals, and no-regret guarantees for online game-theoretic learning achieving sublinear regret growth rates.

Key theorems established fundamental relationships between competing objectives. The Robustness-Accuracy Trade-off Theorem formally characterized the tension between standard accuracy on benign samples and robust accuracy under adversarial perturbations, demonstrating that improving robustness typically degrades standard accuracy with magnitude bounded by local Lipschitz constants. The Privacy-Utility Trade-off analysis quantified utility degradation under differential privacy mechanisms, establishing that strong privacy guarantees require noise levels degrading detection performance with degradation scaling as $O(\sqrt{\frac{\log(1/\delta)}{\epsilon}})$. The Uncertainty Decomposition Theorem formalized separation of total predictive uncertainty into epistemic model uncertainty reducible through additional training data and aleatoric data uncertainty representing inherent noise. These theoretical results provide rigorous foundations supporting principled deployment in security-critical contexts.

\subsection*{Algorithmic Innovations}

Seven novel methods addressed specific sub-problems while contributing to the unified solution framework.

\textbf{Continuous-Time Temporal Graph Neural Networks} introduced the first application of graph neural ordinary differential equations to network security, enabling principled modeling of irregular temporal dynamics through continuous state evolution. The adjoint sensitivity method enabled memory-efficient training through backward integration computing gradients without storing intermediate states. Temporal adaptive batch normalization resolved incompatibilities between continuous dynamics and normalization layers through learned temporal adaptation of scale and shift parameters. Multi-scale temporal modeling captured attack patterns unfolding across time scales from microseconds to months through parallel ODE branches with distinct time constants. The architecture achieved 98.3\% accuracy on encrypted traffic analysis while processing 8.7 million events per second with 47 millisecond median latency enabling real-time enterprise deployment.

\textbf{Triple-Embedding Temporal Graph Neural Networks} pioneered multi-granularity graph representations jointly modeling service-level, trace-level, and node-level security behaviors in microservices environments. Specialized encoders maintained hierarchical embeddings updated through dual temporal-spatial attention mechanisms. Cross-granularity message passing enabled information flow across abstraction levels through bipartite graph attention. Elastic weight consolidation prevented catastrophic forgetting under continuous topology evolution maintaining 94.7\% accuracy despite hundreds of deployment events. The architecture detected 96.8\% of microservices intrusions outperforming single-granularity baselines by 8.3\% through complementary representations capturing multi-level attack characteristics.

\textbf{Federated Learning with Large Language Models} integrated Byzantine-robust optimization, differentially private optimal transport, and semantic reasoning for collaborative threat detection under adversarial conditions. Coordinate-wise trimmed mean aggregation filtered corrupted gradients from malicious participants achieving 87.1\% accuracy with 40\% Byzantine adversaries compared to 68.3\% for standard approaches. Sinkhorn optimal transport aligned heterogeneous distributions across organizations minimizing Wasserstein distance through entropic regularization. Large language model integration enabled 87.6\% F1-score on zero-shot detection of novel attack types through semantic understanding of threat principles. Differential privacy with $(\epsilon = 0.85, \delta = 10^{-5})$ preserved privacy while degrading accuracy by only 3.1\% establishing favorable privacy-utility trade-offs.

\textbf{Post-Quantum Intrusion Detection} developed specialized analysis capabilities for lattice-based, code-based, and hash-based cryptographic protocols addressing emerging threats from quantum computing. Protocol-specific feature extraction captured timing distributions, structural properties, and resource consumption patterns characteristic of post-quantum implementations. Multi-task learning predicted both attack type and target protocol through shared representations enabling transfer learning across cryptographic families. The system detected 96.2\% of attacks against CRYSTALS-Kyber, 93.8\% against Classic McEliece, and 94.7\% against Sphincs+, establishing foundational capabilities for quantum-safe network security. Downgrade attack detection achieved 98.3\% accuracy preventing forced fallback to vulnerable classical protocols.

\textbf{MambaShield} integrated selective state space models with progressive adversarial robustness distillation and PAC-Bayesian certification providing formal generalization guarantees. Input-dependent state transition matrices enabled adaptive filtering of relevant temporal information while maintaining linear computational complexity. Progressive distillation from ensemble teachers achieved 91.4\% accuracy under 25\% label poisoning compared to 73.8\% for standard training through gradual curriculum learning. PAC-Bayesian analysis provided finite-sample risk bounds with 91.7\% empirical coverage probability enabling reliable confidence intervals for security-critical decisions. The architecture combined efficiency of state space models with robustness essential for adversarial environments.

\textbf{Stochastic Transformer} introduced variational Bayesian attention mechanisms enabling uncertainty quantification while maintaining computational tractability. Learned Gaussian distributions over query, key, and value projections propagated uncertainty through attention computation. Monte Carlo sampling with 50 forward passes approximated expectations decomposing total uncertainty into epistemic and aleatoric components. Stochastic adversarial training sampling perturbation parameters from learned distributions achieved 89.6\% robust accuracy compared to 76.3\% for standard approaches. Multi-objective optimization balanced task performance, KL regularization, adversarial robustness, calibration quality, and parameter regularization achieving Expected Calibration Error of 0.078 compared to 0.192 for baseline transformers. Cross-dataset evaluation demonstrated 94.2\% average accuracy with minimal domain-specific tuning.

\textbf{Game-Theoretic Evasion Resistance} characterized equilibrium strategies under adaptive adversaries through Stackelberg game formulations and online learning. Nash equilibrium mixed strategies achieved 94.7\% accuracy against equilibrium evasion strategies compared to 87.6\% for naive approaches. Multiplicative weights online learning maintained sublinear regret $O(\sqrt{T})$ enabling adaptation to evolving threats without complete knowledge of adversary capabilities. Integration with Bayesian uncertainty quantification enabled risk-aware policy selection rejecting predictions with high epistemic uncertainty for human analyst review, reducing investigation time by 43\% while maintaining 97.2\% recall.

\subsection*{Empirical Validation}

Comprehensive experimental evaluation across three major benchmark datasets established practical effectiveness and operational viability. The unified 23-metric evaluation framework characterized performance across detection accuracy, calibration quality, adversarial robustness, computational efficiency, and privacy preservation providing nuanced assessment beyond single-dimensional accuracy.

On CIC-IoT-2023 with 46.6 million flow records, the integrated approach achieved 94.7\% accuracy compared to best baseline at 91.8\% representing 2.9 percentage point improvement. False positive rate of 1.4\% proved acceptable for operational deployment compared to 3.7\% for baselines. Processing throughput of 8.7 million events per second enabled real-time enterprise-scale analysis. Energy consumption of 34 watts per 1000 predictions supported battery-powered IoT edge deployments.

On CSE-CICIDS2018 with 16.2 million enterprise flows, accuracy reached 93.9\% versus 89.7\% for best baseline yielding 4.2 point advantage. The improvement proved especially pronounced for difficult categories including infiltration, botnet detection, and heartbleed exploitation. Cross-dataset generalization maintained 92.7\% accuracy when trained on CIC-IoT-2023 and tested on CSE-CICIDS2018 without retraining demonstrating 1.2\% degradation compared to 8.3\% for baselines.

On UNSW-NB15 with 2.5 million hybrid real-synthetic records, performance achieved 93.8\% accuracy compared to 88.2\% baseline representing 5.6 point gain. Consistent improvements across all nine attack categories confirmed broad applicability. The continuous-time modeling effectively handled irregular inter-arrival times characteristic of this dataset achieving superior temporal alignment.

Systematic ablation studies quantified contributions of individual components confirming architectural necessity. Removing continuous-time dynamics degraded accuracy by 6.8\%, disabling multi-granularity embeddings reduced performance by 8.3\%, eliminating Byzantine-robust aggregation collapsed accuracy by 21.4\% under adversarial conditions, removing optimal transport alignment decreased cross-domain performance by 15.8\%, disabling LLM integration reduced zero-shot F1-score by 53.4\%, removing progressive distillation decreased poisoning robustness by 9.7\%, and replacing variational attention eliminated calibration quality degrading ECE by 114 points.

Cross-domain transfer learning validated general applicability beyond network security. Application to speech phoneme boundary detection achieved 94.2\% F1-score demonstrating utility of continuous-discrete temporal modeling for irregular events. Healthcare adverse event prediction reached 89.7\% accuracy with uncertainty quantification supporting clinical decision making. The successful transfer confirmed architectural generality and broad utility across domains sharing characteristics of irregular temporal dynamics and uncertainty quantification requirements.

\section*{Implications for Network Security Practice}

The research contributions have important implications for contemporary network security operations addressing practical deployment challenges.

\subsection*{Real-Time Threat Detection at Scale}

The demonstrated processing throughput of 8.7 million events per second with sub-100 millisecond latency enables real-time analysis of enterprise network traffic without introducing unacceptable delays. Previous deep learning approaches required offline batch processing limiting responsiveness to active attacks. The continuous-time formulation with efficient ODE integration achieves favorable computational complexity enabling streaming analysis. Organizations can deploy the technology without significant infrastructure investments beyond typical GPU resources already available in many security operations centers.

\subsection*{Privacy-Preserving Collaborative Defense}

The Byzantine-robust federated learning with differential privacy enables practical multi-organization threat intelligence sharing addressing critical barrier to collaborative defense. Regulatory constraints including GDPR prevent centralized data aggregation making traditional machine learning approaches infeasible. The demonstrated $(\epsilon = 0.85, \delta = 10^{-5})$ differential privacy with only 3.1\% accuracy degradation establishes that strong privacy guarantees need not sacrifice detection effectiveness. Information sharing consortiums in finance, healthcare, and critical infrastructure sectors can leverage federated approaches to improve collective security without exposing sensitive network topology or traffic patterns.

\subsection*{Zero-Shot Detection of Novel Threats}

Integration of large language models achieving 87.6\% F1-score on zero-shot detection addresses fundamental limitation of supervised learning approaches requiring labeled examples of all attack types. New vulnerabilities, exploits, and attack techniques emerge continuously faster than security teams can generate comprehensive labeled training data. Semantic understanding through natural language reasoning about attack principles enables proactive detection of previously unseen threats. Security operations centers can maintain effectiveness against evolving attack landscapes without constant model retraining as new threat intelligence arrives.

\subsection*{Quantum-Safe Security Monitoring}

The specialized post-quantum intrusion detection capabilities address imminent transition to quantum-resistant cryptographic protocols. Organizations deploying CRYSTALS-Kyber, Dilithium, Sphincs+, and other NIST standardized post-quantum algorithms require adapted monitoring detecting attacks exploiting new implementation vulnerabilities. Timing side-channels against lattice-based schemes present particularly concerning attack surface requiring millisecond-granularity analysis. The demonstrated 96.2\% detection accuracy on post-quantum attacks establishes viability of adapted intrusion detection for quantum-safe deployments.

\subsection*{Reliable Confidence Estimation for Automated Response}

The calibrated uncertainty quantification with Expected Calibration Error of 0.078 enables trustworthy confidence estimates supporting automated response decisions. Security automation systems require reliable confidence to determine appropriate actions ranging from logging alerts to blocking traffic to isolating compromised systems. Poor calibration causes either excessive false positives triggering unnecessary responses and alert fatigue, or missed attacks due to overconfident incorrect predictions. The demonstrated epistemic-aleatoric uncertainty decomposition provides actionable guidance on when model predictions suffice for automation versus requiring human analyst judgment.

\subsection*{Adaptive Defense Against Strategic Adversaries}

Game-theoretic equilibrium strategies prevent exploitation by adaptive attackers who observe and counter predictable defense patterns. Static detection models become ineffective as sophisticated adversaries develop evasion techniques targeting specific classifiers. The online learning with sublinear regret enables continuous adaptation maintaining effectiveness despite evolving adversary strategies. Integration with uncertainty quantification supports risk-aware policy selection under incomplete information about attacker capabilities and objectives. Security organizations facing advanced persistent threats benefit from strategic reasoning accounting for adversarial adaptation.

\section*{Limitations and Open Challenges}

Despite contributions, several limitations merit acknowledgment and suggest directions for continued research.

\subsection*{Computational Overhead}

Continuous-time neural ODE integration requires adaptive ODE solvers introducing computational overhead compared to discrete architectures. Training time increases by approximately 1.8× despite mathematical elegance of continuous formulations. The adjoint sensitivity method mitigates memory costs but gradient computation remains more expensive than standard backpropagation. Monte Carlo uncertainty quantification through 50 forward passes introduces 8× inference slowdown after batch processing optimization. Practical deployment requires balancing accuracy gains against computational costs, with edge devices facing tighter constraints than datacenter infrastructure.

\subsection*{Hyperparameter Sensitivity}

Federated learning performance exhibits sensitivity to hyperparameter selection including trimming fractions for Byzantine-robust aggregation, learning rates, privacy noise scales, and optimal transport regularization strengths. Cross-validation on federated data proves challenging when participants cannot share raw samples. Automated hyperparameter optimization through grid search or Bayesian optimization incurs significant communication overhead in distributed settings. Practitioners require guidance on hyperparameter selection for specific deployment contexts balancing multiple objectives.

\subsection*{Interpretability Challenges}

Deep continuous models with neural ODEs and attention mechanisms provide limited interpretability compared to rule-based systems or decision trees. Security analysts require explanations of detection decisions to validate alerts, investigate incidents, and refine policies. Attention weight visualization provides partial interpretability but continuous state evolution trajectories resist intuitive explanation. Regulatory compliance requirements in critical sectors mandate explainable AI systems. Future work should develop interpretability techniques specifically adapted to continuous-time architectures including attention visualization, saliency analysis, and counterfactual explanations.

\subsection*{Adversarial Robustness Gaps}

Despite improvements over baselines, robust accuracy of 89.6\% under adversarial attacks leaves 10.4\% vulnerability representing millions of potential evasion opportunities at enterprise scale. Adaptive attacks that observe defense mechanisms and develop customized evasion strategies may achieve higher success rates than evaluated white-box attacks. The robustness-accuracy trade-off theorem establishes fundamental limits suggesting perfect robustness and accuracy may prove unachievable simultaneously for certain data distributions. Continued research on certified defenses with tighter bounds and adaptive training procedures remains critical.

\subsection*{Concept Drift Adaptation}

Evaluation focused primarily on static datasets with fixed data distributions. Operational networks exhibit continuous evolution through infrastructure changes, user behavior shifts, and emerging attack techniques constituting concept drift. Online learning algorithms require careful design balancing plasticity for new patterns with stability preventing catastrophic forgetting of critical historical threats. The elastic weight consolidation approach provides partial solution but more sophisticated continual learning techniques warrant investigation. Automated detection of distribution shift triggering model updates would enhance operational viability.

\subsection*{Multi-Modal Integration}

Current approaches focus primarily on network traffic analysis despite security operations centers having access to diverse data sources including system logs, threat intelligence feeds, vulnerability scans, and endpoint telemetry. Multi-modal fusion correlating complementary information sources could significantly improve detection through cross-validation and enriched context. However, heterogeneous data types with varying semantics, temporal granularities, and reliability characteristics pose integration challenges. Unified representations spanning modalities while preserving modality-specific characteristics require careful architectural design.

\section*{Future Research Directions}

Several promising directions warrant investigation in future research including planned postdoctoral work.

\subsection*{Stochastic Differential Equations for Uncertainty}

Extending deterministic neural ODEs to stochastic differential equations~\cite{li2020scalable} enables principled uncertainty modeling through diffusion processes. The Langevin dynamics~\cite{welling2011bayesian} formulation $\frac{d\mathbf{x}(t)}{dt} = f(\mathbf{x}(t), t) + \sigma(t) \frac{d\mathbf{W}(t)}{dt}$ where $\mathbf{W}(t)$ denotes Wiener process introduces stochasticity capturing epistemic uncertainty directly in continuous dynamics. Fokker-Planck equations~\cite{risken1996fokker} characterize evolution of probability distributions over states enabling analytical uncertainty propagation. Practical implementation requires specialized numerical integrators for stochastic processes and careful analysis of discretization errors. The approach promises tighter integration of continuous-time modeling with Bayesian inference compared to current Monte Carlo sampling.

\subsection*{Unbalanced Optimal Transport}

Classical optimal transport assumes equal total mass in source and target distributions limiting applicability to imbalanced datasets common in security where attacks constitute small minorities. Unbalanced optimal transport~\cite{chizat2018unbalanced,sejourne2019sinkhorn} relaxes mass preservation constraints through regularization terms $\int d\mu_s + \lambda_1 KL(\mu_s|\alpha) + \int d\mu_t + \lambda_2 KL(\mu_t|\beta)$ enabling mass creation and destruction. The extended formulation naturally handles class imbalance during cross-domain adaptation. Entropic unbalanced optimal transport admits efficient Sinkhorn-like algorithms maintaining computational tractability. Investigation of unbalanced transport for security applications with extreme class imbalance ratios exceeding 1000:1 would enhance cross-dataset generalization.

\subsection*{Foundation Models for Security}

Large-scale pre-training on massive security data corpora could yield foundation models~\cite{bommasani2021opportunities} analogous to language model pre-training~\cite{devlin2019bert,brown2020language} on internet text. Masked flow prediction, contrastive learning~\cite{chen2020simple} on similar attack sequences, and denoising autoencoders~\cite{vincent2008extracting} on corrupted traffic could extract general representations transferable across security domains. The Cybersecurity Knowledge Graph integrating threat intelligence, vulnerability databases, attack patterns, and defense strategies provides structured knowledge for grounding. Fine-tuning foundation models on specific deployment contexts with limited labeled data would reduce annotation requirements. However, computational costs of billion-parameter security models and privacy concerns about training data require careful consideration.

\subsection*{Quantum Machine Learning}

Quantum computing threatens current cryptography but also offers computational advantages for machine learning~\cite{biamonte2017quantum}. Quantum kernel methods~\cite{havlicek2019supervised} computing inner products in exponentially large Hilbert spaces enable richer feature representations than classical kernels. Variational quantum circuits~\cite{mcclean2016theory} provide parameterized quantum states optimized through hybrid classical-quantum algorithms. Quantum advantage demonstrations for classification tasks establish practical utility despite current hardware limitations. Investigation of quantum machine learning for intrusion detection could leverage quantum speedups for challenging optimization problems including adversarial training where finding worst-case perturbations requires expensive optimization. Near-term noisy intermediate-scale quantum devices suffice for initial experiments.

\subsection*{Human-AI Collaboration}

Security operations require human analysts reviewing alerts, investigating incidents, and making critical response decisions. Uncertainty-aware systems providing calibrated confidence estimates enable effective human-AI collaboration~\cite{bansal2021does} with automated handling of high-confidence predictions and human review of uncertain cases. Active learning algorithms~\cite{settles2009active} selecting informative samples for human labeling improve model quality efficiently. Interactive machine learning~\cite{amershi2014power} incorporating analyst feedback in real-time accelerates adaptation to emerging threats. Explainable AI techniques~\cite{ribeiro2016should,lundberg2017unified} visualizing model reasoning support analyst understanding and trust. Designing human-AI workflows optimizing complementary strengths of automation and human expertise remains crucial challenge. Cognitive psychology research on human trust calibration and decision making under uncertainty informs effective interface design.

\subsection*{Adversarial Machine Learning Theory}

Despite empirical progress, theoretical understanding of adversarial robustness remains incomplete~\cite{schmidt2018adversarially,zhang2019theoretically}. Fundamental questions include: Do robust classifiers exist for given data distributions? What is the optimal achievable robust accuracy? How does sample complexity scale with robustness requirements? Can we characterize intrinsic data properties determining robustness limits? Connections to computational learning theory~\cite{valiant1984theory}, statistical physics~\cite{shalev2014understanding}, and algorithmic game theory~\cite{nisan2007algorithmic} provide analytical tools. Establishing impossibility results clarifying fundamental barriers would guide practical expectations. Conversely, proving achievability theorems with constructive algorithms would establish attainable goals. Advances in adversarial ML theory would ground empirical security research in rigorous foundations.

\subsection*{Continuous-Time Explainability}

Explaining predictions from continuous-time models requires new interpretability techniques adapted to infinite-dimensional trajectory optimization. Sensitivity analysis~\cite{saltelli2008global} through adjoint equations computes feature importance gradients efficiently. ODE vector field visualization~\cite{chen2018neural} reveals learned dynamics geometry. Trajectory clustering identifies canonical temporal patterns. Attention flow tracking through continuous time shows information aggregation evolution. Counterfactual trajectory generation~\cite{wachter2017counterfactual} demonstrates minimal input modifications changing predictions. Integration with domain knowledge about attack stages~\cite{hutchins2011intelligence} and security principles grounds explanations in actionable insights. Developing comprehensive explainability framework for continuous-time security models would support analyst trust and regulatory compliance.

\section*{Broader Impact}

The research advances fundamental capabilities for protecting networked systems against evolving cyber threats with implications extending beyond technical contributions.

\subsection*{Critical Infrastructure Protection}

Electrical grids, water treatment facilities, transportation systems, and healthcare infrastructure increasingly rely on networked digital control systems vulnerable to cyber attacks~\cite{colbert2016critical}. State-sponsored adversaries and criminal organizations target critical infrastructure~\cite{lewis2014critical} seeking disruption or sabotage. The demonstrated real-time detection capabilities with formal robustness guarantees support deployment in safety-critical contexts where failures endanger public safety. Privacy-preserving federated learning enables information sharing among critical infrastructure operators maintaining operational security while improving collective defense posture. Quantum-safe monitoring addresses long-term threats as quantum computers mature threatening current cryptographic protections.

\subsection*{National Security}

Nation-state cyber operations represent strategic threats to economic competitiveness, military effectiveness, and political stability~\cite{rid2012cyber}. Advanced persistent threats~\cite{tankard2011advanced} employ sophisticated multi-stage attacks requiring extended dwell times for reconnaissance, lateral movement, and exfiltration. The multi-granularity temporal modeling detecting coordinated activities across abstraction levels addresses advanced threat detection. Game-theoretic equilibrium strategies account for strategic adversarial adaptation by sophisticated actors. Zero-shot detection through semantic reasoning enables identification of novel attack techniques before comprehensive labeled training data becomes available. Enhanced defensive capabilities strengthen deterrence by increasing attack costs and reducing success probabilities.

\subsection*{Economic Security}

Cybercrime costs global economy hundreds of billions annually~\cite{anderson2019measuring} through direct theft, ransomware~\cite{cartwright2019cyber}, fraud, and operational disruptions. Small and medium enterprises lack security expertise and resources of large corporations making them attractive targets. The computational efficiency enabling edge deployment and automated operation reduces barriers to adoption of effective defenses. Privacy-preserving collaborative detection allows industry sectors to share threat intelligence without exposing competitive information. Reliable automated response based on calibrated confidence enables security automation for organizations without dedicated security operations centers. Improved cybersecurity supports economic growth by reducing digital risks and enabling innovation.

\subsection*{Digital Rights and Privacy}

Mass surveillance~\cite{lyon2015surveillance} and data breaches threaten individual privacy and civil liberties in increasingly digital societies. Differential privacy mechanisms~\cite{dwork2006calibrating,dwork2014algorithmic} providing formal guarantees enable security monitoring while protecting sensitive information about communication patterns, behavioral patterns, and personal data. The favorable privacy-utility trade-offs demonstrate that security and privacy need not be opposing objectives but can be achieved simultaneously through careful system design. Federated approaches distributing computation avoid centralized data aggregation reducing risks of bulk surveillance. Enhanced security technologies protect citizens against criminal exploitation while respecting privacy rights and democratic values.

\subsection*{Ethical Considerations}

Machine learning for security raises ethical considerations~\cite{jobin2019global} requiring thoughtful governance. Adversarial robustness techniques could enable censorship-resistant communications circumventing authoritarian surveillance if repurposed by activists. Dual-use concerns~\cite{brundage2018malicious} mandate responsible disclosure and access controls preventing malicious applications. Bias in training data~\cite{barocas2016big} may lead to discriminatory detection patterns unfairly targeting specific demographic groups or organizations. Automated response systems making blocking decisions without human oversight risk denial of service to legitimate users. Transparency requirements balanced against security imperatives create tensions between explainability and protecting defense capabilities. Ongoing ethical analysis incorporating diverse stakeholder perspectives should guide responsible deployment supporting security objectives while respecting human rights and societal values.

\section*{Concluding Remarks}

This dissertation established theoretical foundations and practical implementations for adversarially resilient artificial intelligence in hybrid intrusion detection and prevention systems. The unified framework integrating continuous-time graph neural networks, optimal transport theory, Bayesian inference, and game-theoretic modeling provides comprehensive solution to fundamental challenges of adversarial robustness, privacy preservation, cross-domain generalization, and computational efficiency. Seven novel methods addressed specific sub-problems while contributing to integrated defense capabilities validated through extensive experimentation demonstrating state-of-the-art performance across multiple benchmark datasets.

The research advances network security practice by enabling real-time threat detection at enterprise scale, privacy-preserving collaborative defense, zero-shot detection of novel attacks, quantum-safe security monitoring, reliable confidence estimation for automation, and adaptive strategies against sophisticated adversaries. While limitations remain including computational overhead, hyperparameter sensitivity, interpretability challenges, and adversarial robustness gaps, the demonstrated capabilities represent significant progress toward trustworthy AI for security-critical applications.

Future research directions including stochastic differential equations, unbalanced optimal transport, security foundation models, quantum machine learning, human-AI collaboration, adversarial theory, and continuous-time explainability promise continued advancement. The broader implications for critical infrastructure protection, national security, economic security, digital rights, and ethical governance underscore importance of rigorous research grounded in mathematical foundations while addressing practical deployment requirements and societal considerations.

As networks grow increasingly complex and adversaries more sophisticated, adversarially resilient artificial intelligence will prove essential for maintaining security, privacy, and trust in digital systems underpinning modern society. This dissertation contributes both theoretical insights and practical tools advancing that critical objective while establishing foundations for future innovations in this vital domain.
